\subsubsection{Desplazamiento de prioridades y persistencia del temor}
Entre 2019 y 2021, la economía como principal problema aumenta 14,7 puntos,
mientras seguridad disminuye 33,2 puntos (figura \ref{fig:prioridades}). En
2021, COVID-19 representa 22,0\,\%. La incorporación de esa categoría cambia
el conjunto de respuestas: las variaciones de seguridad y economía no son
una redistribución exclusiva entre ambas. Tampoco miden la intensidad de
cada preocupación, pues elegir un problema principal puede coexistir con
preocupaciones importantes en otros ámbitos.

\begin{figure}[htbp]\centering
\begin{tikzpicture}\begin{axis}[width=0.85\textwidth,height=5.7cm,
 xmin=2016.8,xmax=2021.2,ymin=0,ymax=80,xtick={2017,2018,2019,2021},
 scaled x ticks=false,/pgf/number format/1000 sep={},ylabel={Porcentaje},
 grid=major,grid style={gray!15},tick label style={font=\small},
 legend style={at={(0.5,-0.18)},anchor=north,legend columns=3,column sep=6pt,draw=none,font=\small}]
\addplot[petroleo,thick,mark=*] table[x=anio,y=seguridad,col sep=comma]{data/infosegura_preocupaciones.csv};
\addlegendentry{Seguridad}
\addplot[ocre,thick,mark=square*] table[x=anio,y=economia,col sep=comma]{data/infosegura_preocupaciones.csv};
\addlegendentry{Economía}
\addplot[gris,only marks,mark=triangle*,mark size=3pt] coordinates {(2021,22.02)};
\addlegendentry{COVID-19}
\end{axis}\end{tikzpicture}
\caption{Evolución de las prioridades declaradas: principal problema del país}\label{fig:prioridades}
\fuente{PNUD/Infosegura (2021). Series seleccionadas, no distribución completa de respuestas. COVID-19 solo se observa como categoría en 2021; los años anteriores no se codifican como cero. Sin dato de 2020 ni IC propios.}
\end{figure}

La percepción de inseguridad también disminuye, pero conserva diferencias
marcadas por espacio (figura \ref{fig:percepcion}). Entre 2019 y 2021 cae
10,7 puntos en transporte público y 12,5 en comunidad. En 2021, las cifras
reportadas son 68,1\,\% y 19,9\,\%, respectivamente. Su distancia descriptiva
de 48,2 puntos muestra que un mismo diagnóstico nacional requiere observar
lugares de movilidad y de residencia. La fuente no detalla las bases de
cada espacio; la distancia no se interpreta como una comparación de riesgo
individual ajustada por exposición.

\begin{figure}[htbp]\centering
\begin{tikzpicture}\begin{axis}[width=0.85\textwidth,height=5.7cm,
 xmin=2016.8,xmax=2021.2,ymin=0,ymax=100,xtick={2017,2018,2019,2021},
 scaled x ticks=false,/pgf/number format/1000 sep={},ylabel={Porcentaje reportado},
 grid=major,grid style={gray!15},tick label style={font=\small},
 legend style={at={(0.5,-0.18)},anchor=north,legend columns=2,column sep=6pt,draw=none,font=\small}]
\addplot[petroleo,thick,mark=*] table[x=anio,y=transporte_publico,col sep=comma]{data/infosegura_percepcion.csv};
\addlegendentry{Transporte público}
\addplot[ocre,thick,mark=square*] table[x=anio,y=comunidad,col sep=comma]{data/infosegura_percepcion.csv};
\addlegendentry{Comunidad}
\end{axis}\end{tikzpicture}
\caption{Percepción de inseguridad en transporte y comunidad}\label{fig:percepcion}
\fuente{PNUD/Infosegura (2021). La base de personas de 18+ del indicador de uso de espacios no se traslada a estas series de percepción. Percepción distinta de victimización; segmentos entre años observados, sin imputar 2020.}
\end{figure}

La combinación de menor victimización declarada y persistencia del temor en
transporte sugiere evaluar varias dimensiones de seguridad a la vez. Para
las familias, el acceso a escuela, empleo y servicios requiere movilidad;
el temor puede ser relevante incluso sin una experiencia delictiva reciente.
Esta es una interpretación para orientar preguntas de investigación, no una
estimación del efecto del transporte sobre disciplina o primera unión. Se
necesitaría medir uso efectivo, frecuencia de desplazamiento y horarios para
examinar esa relación.

\subsubsection{Privaciones: cambios heterogéneos entre 2019 y 2020}
La figura \ref{fig:privaciones} examina los veinte cambios de privación disponibles.
Siete indicadores aumentan, once disminuyen y dos quedan sin cambio en la
precisión publicada. El aumento de privación por desempleo (+10,4 puntos)
destaca frente a otras variaciones positivas de hasta 1,3 puntos. Al mismo
tiempo, disminuyen restricciones por inseguridad ($-6{,}8$), falta de acceso
a salud ($-3{,}0$) e incidencia de crimen y delito ($-2{,}4$). El conjunto
documenta movimientos de distinto signo, relevantes para una interpretación
multidimensional del bienestar.

\begin{figure}[!htbp]\centering
\begin{tikzpicture}\begin{axis}[width=0.76\textwidth,height=11.0cm,xbar,
 bar width=7pt,xmin=-8,xmax=12,xtick={-8,-4,0,4,8,12},
 ymin=-0.7,ymax=19.7,ytick={0,...,19},
 yticklabels={Restricciones por inseguridad,Sin acceso a salud,Crimen y delito,
 Hacinamiento,Tenencia del terreno,Cuidado temprano inadecuado,
 Sin seguridad social,Sin saneamiento,Inasistencia escolar,
 Subempleo e inestabilidad,Trabajo infantil,Rezago educativo,
 Sin agua potable,Pared y piso inadecuados,Sin espacios de esparcimiento,
 Baja educación de adultos,Techo inadecuado,Inseguridad alimentaria,
 Daños y riesgos ambientales,Desempleo},
 xlabel={Variación 2019--2020 (puntos porcentuales)},
 axis x line*=bottom,axis y line*=left,xmajorgrids=true,grid style={gray!15},
 tick label style={font=\scriptsize},nodes near coords,
 nodes near coords style={font=\scriptsize,/pgf/number format/fixed,/pgf/number format/precision=1}]
\addplot[fill=petroleo,draw=none] table[x=cambio,y=pos,col sep=comma]{data/derivada_privaciones.csv};
\draw[gris,thick] (axis cs:0,-0.7)--(axis cs:0,19.7);
\end{axis}\end{tikzpicture}
\caption{Variación de los indicadores de privación}\label{fig:privaciones}
\fuente{PNUD/Infosegura (2021). Etiquetas abreviadas; textos originales en CSV. Cambios de indicadores de privación, no sus niveles. El componente desempleo no es la tasa de desempleo. Cero es ausencia de cambio al redondeo disponible.}
\end{figure}

La magnitud de desempleo invita a examinar apoyos de ingreso y cuidado junto
con servicios de protección. Sin embargo, no se suman las veinte variaciones
para obtener una variación de pobreza multidimensional: las privaciones
pueden coexistir en un hogar y el índice depende de ponderaciones, umbral y
distribución conjunta. Una reducción del indicador de inseguridad tampoco
implica que el nivel de privación sea bajo o que haya desaparecido. Sin
niveles iniciales, no pueden compararse niveles finales solo por la altura
de estas barras.

La lectura conjunta con MICS permite formular una agenda concreta: distinguir
estrés material, aceptación de prácticas violentas y exposición efectiva;
observar sus intersecciones con datos compatibles; y evaluar si los servicios
alcanzan a hogares con varias necesidades. Los agregados de 2019--2021 no
explican por sí solos la conducta familiar de 2014. Su contribución al
artículo consiste en precisar las dimensiones que una siguiente medición
debe integrar, manteniendo sus distintos períodos y poblaciones.
